\section{Background}
\label{sec:background}

\subsection{Self-tracking has become an ordinary data practice}

Health and fitness tracking is no longer a niche activity. In Germany,
73\,\% of smartphone users report using health apps to support fitness,
nutrition, or general wellbeing~\parencite{bitkom2025health}. The category is
broad: step and activity counters, sleep and heart-rate monitors, workout and
nutrition apps, and menstrual and fertility trackers. A wearable device may or
may not be involved. What is always involved is an app --- it is where data is
entered, displayed, stored, and transmitted. This is why the unit of analysis
in this thesis is the app and the communication surrounding it rather than the
hardware.

Three kinds of data accumulate. Users supply attributes directly: age, weight,
goals, symptoms, mood, sexual activity. Sensors produce continuous streams:
steps, heart rate, sleep phases, GPS traces. The system adds its own context:
timestamps, device identifiers, usage patterns~\parencite{gabriele2020understanding}.
What distinguishes tracking data from most other consumer data is that it is
longitudinal and involuntary in character --- once the app is installed, the
record keeps growing whether or not the user is thinking about privacy. Under
Article~9(1) of the General Data Protection Regulation, data concerning health
is a special category of personal data; menstrual and fertility data falls
squarely within it~\parencite{gdpr2016}.

\subsection{From data practices to privacy harms}

A \emph{privacy violation} and a \emph{privacy harm} are not the same thing.
Solove's taxonomy classifies the harmful \emph{activities} --- information
collection, processing, dissemination, and
invasion~\parencite{solove2006taxonomy}. Citron and Solove later classify the
\emph{injuries} those activities produce, and it is this second typology that
the present work builds on: physical, economic, reputational, psychological,
autonomy, discrimination, and relationship
harms~\parencite[830--861]{citron2021harms}. Calo had earlier drawn a related
line between the \emph{subjective} harm of perceived unwanted observation and
the \emph{objective} harm of information being used against a
person~\parencite{calo2011boundaries}.

What these typologies do not supply is the route from a system property to an
injury. Knabenhans et al.\ close that gap with \emph{harm
trees}~\parencite{knabenhans2026harms}. A tree links system components ---
information the functionality necessarily leaks, requirements the system
imposes on users --- to adversarial capabilities, and those capabilities to
harms. Their analysis of the European Digital Identity Framework isolates three
capabilities that follow from the core functionality itself: distinguishing
between users, tracking users, and inferring attributes about them. Crucially,
they separate pathways that are inherent to the functionality, pathways that
arise from deploying it in a digital context, and pathways that come from
concrete implementation choices --- and show that privacy-enhancing
technologies (PETs) can only close the last of the three. The same three
capabilities describe a health tracker well, which makes the model transferable
to this domain even though the original object of study is different.

That inference capability is not hypothetical here. Salehzadeh Niksirat et al.\
reviewed 236 studies of wearable activity trackers and report that tracker data
can be mined to infer private attributes, personality traits among
them~\parencite{niksirat2024wearable}.

\subsection{Modern risk is inference-centric, and so is PET marketing}

Martin, Nissenbaum, and Shmatikov examine three principles that underlie the
privacy techniques large technology companies have recently deployed:
limiting third-party access to personal data, minimising the use and retention
of raw data in favour of inferences drawn from it, and keeping personal data on
the user's device~\parencite{martin2025cookies}. Their objection is not that the
techniques fail on their own terms. It is that a technique can satisfy all
three principles while the privacy-violating behaviour continues untouched,
because the sensitive object is the inference, not the raw record. Their
vignette studies show that users' expectations do not follow the first
party\,/\,third party or raw data\,/\,inference distinctions that the principles
are built on.

For a health tracker this is the central point. Cycle length and symptom entries
support pregnancy inference; location and activity traces support inferences
about routine, workplace, and health status. Prior work on period trackers has
documented both the technical basis for such inferences and the concrete fears
users attach to them~\parencite{cao2024deleted,shipp2020period}.

\subsection{Where privacy is communicated}

Providers communicate privacy on several surfaces at once, and the surfaces
differ in audience, length, and legal weight:

\begin{enumerate}[leftmargin=*,itemsep=2pt]
  \item the \textbf{privacy policy}, the only surface with a direct legal
        mandate;
  \item the \textbf{store listing} --- the description text plus the structured
        disclosure that each platform requires (Google Play's Data Safety
        section, Apple's Privacy Nutrition Labels);
  \item the \textbf{provider's website and marketing material}, including
        advertisements and company privacy blog posts;
  \item the \textbf{app interface itself} --- onboarding screens, consent and
        permission dialogues, privacy settings, and the paths that lead to
        export or deletion.
\end{enumerate}

Each surface has known weaknesses. Policies are long enough that reading them is
economically irrational~\parencite{mcdonald2008cost}. Structured store
disclosures are self-reported and frequently wrong: a dynamic study of almost
5{,}000 Android apps found that 81\,\% misrepresented their data collection or
sharing practices in the Data Safety
section~\parencite{arkalakis2024abandon}, and a large-scale longitudinal study of
Apple's labels documents how unevenly they are created and
maintained~\parencite{li2022ios}.

\subsection{What the law asks providers to say}
\label{sec:law}

The transparency regime is built around describing \emph{what will be done with
data}. Articles 13 and 14 GDPR enumerate the controller's identity, the purposes
and legal basis of processing, recipients, third-country transfers, retention
periods, and data-subject rights~\parencite{gdpr2016}. Consequences appear in
exactly one narrow place: for automated decision-making within the meaning of
Article 22(1) and (4), Article 13(2)(f) requires \enquote{meaningful information
about the logic involved, as well as the significance and the envisaged
consequences of such processing for the data subject}. Breaches are addressed
reactively --- Article 34 obliges the controller to communicate a breach to
affected individuals, but only once it has occurred and only where it is likely
to result in a high risk to their rights and freedoms.

Nothing in the regime requires a provider to state what can be \emph{derived}
from the data it collects, or what would \emph{follow for the user} if that data
were exposed. This is the structural reason to expect silence on inference,
breach consequences, and incident handling --- and the reason that silence is
worth measuring rather than assuming.

\subsection{What prior work establishes, and where it stops}

Three strands of work bear on the question, and each stops short of it.

\paragraph{Stated practice versus actual practice.}
A substantial measurement literature compares what apps say with what they do.
Huckvale et al.\ found that 29 of 36 depression and smoking-cessation apps
transmitted data to services run by Google or Facebook while only 12 disclosed
this in a policy~\parencite{huckvale2019assessment}. Shipp and Blasco
systematically analysed the privacy policies of 30 Android menstrual
apps~\parencite{shipp2020period}. Malki et al.\ combined a thematic analysis of
Data Safety sections and privacy policies with a privacy-focused usability
inspection across 20 popular female mHealth apps and found inconsistencies
between policy text and in-app features, flawed consent and deletion mechanisms,
and covert collection of sensitive data~\parencite{malki2024exploring}. Hudig and
Singh analysed the network transmissions of eight fertility tracker products and
found data flowing to endpoints associated with organisations across several
countries~\parencite{hudig2025intimate}. At scale, MAPS audited
compliance across roughly a million apps~\parencite{zimmeck2019maps}. All of
these ask whether the disclosure is \emph{true}. None asks which \emph{harms} it
names.

\paragraph{What users know and expect.}
Gabriele and Chiasson surveyed 212 fitness tracker users: participants were
reasonably confident about \emph{what} data was collected but unsure \emph{how}
it was used, and when presented with scenarios carrying security and privacy
implications they judged them plausible yet
unlikely~\parencite{gabriele2020understanding}. Cao et al.\ surveyed 183 US
women after the overturn of Roe v.\ Wade and found high concern about period
tracker data practices combined with a feeling of being uninformed and powerless
about mitigation~\parencite{cao2024deleted}. These studies measure the belief;
they do not analyse the material that shapes it.

\paragraph{Methods for reading this material.}
Privacy policies have been a corpus-analysis object for a decade. The OPP-115
corpus established an annotation scheme for data practices in
policies~\parencite{wilson2016creation}; Polisis showed that those categories can
be predicted automatically~\parencite{harkous2018polisis}; and large language
models have since been evaluated for policy analysis at
scale~\parencite{rodriguez2024llm}. On the interface side, the cognitive
walkthrough offers a structured, task-based inspection
procedure~\parencite{wharton1994cognitive}, complementing heuristic
evaluation~\parencite{nielsen1994heuristic}. The annotation schemes these methods
carry, however, are schemes of \emph{data practices} --- collection, sharing,
retention, choice. None of them has a category for harm.
